\section{Освещение сцены свечением сияния}
\label{sec:aurora-lighting}

Освещённость, создаваемая элементом источника, пропорциональна косинусу угла падения и обратно пропорциональна квадрату расстояния до него~\cite[с.~381, 402]{rogers1989}, поэтому освещённость точки $x$ свечением среды равна
\begin{equation}
	E(x) = \int_V \varepsilon \rho(y)\, C(y)\, \frac{\max(0,\ \vec{N} \cdot \vec{L}_y)}{|x - y|^2}\, S(x, y)\, dV,
	\label{eq:aurora-light}
\end{equation}
где $\vec{L}_y$ "--- единичный вектор из $x$ в точку среды $y$, $S(x, y) \in \{0, 1\}$ "--- признак видимости точки $y$ из $x$ (тень).

\subsection{Интегрирование по узлам сетки}

Каждый узел сетки рассматривается как точечный источник, и интеграл~\eqref{eq:aurora-light} заменяется суммой по узлам. Метод точен, но требует до $8{,}4 \cdot 10^6$ слагаемых на точку поверхности и теневого буфера на каждый узел.

\subsection{Набор точечных источников}

Основание каждого занавеса делится на участки, и свечение участка $V_j$ заменяется точечным источником с интенсивностью $I_j$, расположенным в центре яркости участка $y_j$:
\begin{equation}
	I_j = \int_{V_j} \varepsilon \rho\, dV, \qquad y_j = \frac{1}{I_j} \int_{V_j} y\, \varepsilon \rho(y)\, dV.
	\label{eq:point-source}
\end{equation}
Тогда интеграл~\eqref{eq:aurora-light} заменяется суммой по $m$ источникам:
\begin{equation}
	E(x) \approx \sum_{j=1}^{m} I_j\, C_j\, \frac{\max(0,\ \vec{N} \cdot \vec{L}_j)}{|x - y_j|^2}\, S_j(x),
	\label{eq:point-sum}
\end{equation}
где $C_j$ "--- средний цвет участка, $\vec{L}_j$ "--- единичный вектор из $x$ в $y_j$, а признак $S_j$ вычисляется по своему теневому буферу~\eqref{eq:shadow-test}. Погрешность тем меньше, чем меньше участок по сравнению с расстоянием до освещаемой точки.

\subsection{Выбор метода освещения сцены сиянием}

Оценки в таблице~\ref{tab:aurora-light} получены для $p \approx 2{,}1 \cdot 10^6$ пикселей и $m = 16$ источников (по два на занавес).

\begin{table}[H]
	\caption{Сравнение методов освещения сцены свечением сияния}
	\label{tab:aurora-light}
	\begin{tabular}{|L{0.44}|C{0.27}|C{0.27}|}
		\hline
		Критерий & Интегрирование по узлам & Набор точечных источников \\
		\hline
		Слагаемых на кадр & $\approx 1{,}8 \cdot 10^{13}$ & $\approx 3{,}4 \cdot 10^7$ \\
		\hline
		Теневых буферов & по числу узлов & $m$ \\
		\hline
		Точность & точная & приближённая \\
		\hline
	\end{tabular}
\end{table}

Интегрирование по узлам требует примерно в $5 \cdot 10^5$ раз больше слагаемых, поэтому выбран набор точечных источников.